\subsection{\texorpdfstring{THE VECTOR SPACE \(R^{n}\)}{THE VECTOR SPACE R TO THE N}}

Since \(\mathbf{R}\) is a field, we can define on \(\mathbf{R}^n\) a vector space structure over the field \(\mathbf{R}\) , the product \(tx\) of a scalar \(t \in \mathbf{R}\) and a point (or vector) \(\boldsymbol{x} = (x_i)\) of \(\mathbf{R}^n\) being the point \((tx_i)\) . Note that the homothety \((t, x) \to tx\) is continuous on \(\mathbf{R} \times \mathbf{R}^n\) . If \(\boldsymbol{e}_i\) denotes the vector of \(\mathbf{R}^n\) all of whose coordinates are zero, except for that of index \(i\) , which is equal to \(1\) , then the \(\boldsymbol{e}_i\) form a basis of the vector space \(\mathbf{R}^n\) , called the canonical basis of this space. Every vector \(\boldsymbol{x} = (x_i) \in \mathbf{R}^n\) can be written as \(\boldsymbol{x} = \sum_{i=1}^{n} x_i \boldsymbol{e}_i\) , and the relation \(\sum_{i=1}^{n} t_i \boldsymbol{e}_i = 0\) implies that \(t_i = 0\) for \(1 \leqslant i \leqslant n\) .

The vector space \(\mathbf{R}^n\) is therefore of dimension \(n\) over the field \(\mathbf{R}\) , in the sense of algebra (Algebra, Chapter II, § 7, no. 2); hence its name of \(n\) -dimensional real number space.

Let \(f\) be an affine mapping of the vector space \(\mathbf{R}^n\) into the vector space \(\mathbf{R}^m\) ( \(m\) and \(n\) being integers \(>0\) ). If we put \(g(\boldsymbol{x}) = f(\boldsymbol{x}) - f(0)\) , \(g\) is a linear mapping of \(\mathbf{R}^n\) into \(\mathbf{R}^m\) . Let \(a_{ij}\) ( \(1 \leqslant j \leqslant m\) ) be the coordinates of \(g(e_i)\) in \(\mathbf{R}^m\) and let \(b_j\) ( \(1 \leqslant j \leqslant m\) ) be those of \(f(0)\) ; if \(x_i\) ( \(1 \leqslant i \leqslant n\) ) is the \(i\) th coordinate of \(\boldsymbol{x} \in \mathbf{R}^n\) and if \(y_j\) is the \(j\) th coordinate of \(y = f(\boldsymbol{x})\) , we have

\[
y _ {j} = \sum_ {i = 1} ^ {n} a _ {i j} x _ {i} + b _ {j} \quad (1 \leqslant j \leqslant m).
\]

Since every linear polynomial in \(x_1, x_2, \ldots, x_n\) is uniformly continuous on \(\mathbf{R}^n\) , it follows that every affine mapping of \(\mathbf{R}^n\) into \(\mathbf{R}^m\) is uniformly continuous on \(\mathbf{R}^n\) (Chapter II, § 2, no. 6, Proposition 7).

In particular, we know that every affine mapping of \(R^{n}\) onto itself is bijective and that its inverse is again an affine mapping; hence every affine mapping of \(R^{n}\) onto itself is a homeomorphism (and an automorphism of the uniform structure of \(R^{n}\) ).

Let \((\pmb{a}_i)_{1\leqslant i\leqslant n}\) be a free system of \(n\) vectors of \(\mathbf{R}^n\) {[}in other words a basis of the vector space \(\mathbf{R}^n\) {]}; if \(\pmb{b}\) is any point of \(\mathbf{R}^n\) , the set \(\mathbf{P}\) of points \(\pmb{x} = \pmb{b} + \sum_{i=1}^{n} u_i a_i\) such that \(-1 \leqslant u_i \leqslant 1\) for \(1 \leqslant i \leqslant n\) is a compact neighbourhood of \(\pmb{b}\) ; for there exists a bijective affine mapping \(f\) of \(\mathbf{R}^n\) onto itself such that \(f(\pmb{b}) = 0\) , \(f(\pmb{b} + a_i) = e_i\) for \(1 \leqslant i \leqslant n\) ; and \(f(\mathbf{P})\) is the cube which is the product of the \(n\) intervals \([-1, +1]\) in the \(n\) factor spaces. \(\mathbf{P}\) is called the closed parallelotope with centre \(\pmb{b}\)

and basis vectors \(a_{i}\) . The interior of P consists of the points \(b + \sum_{i=1}^{n} u_{i} a_{i}\) such that \(-1 < u_{i} < 1\) for \(1 \leqslant i \leqslant n\) ; it is called the open parallelotope with centre b and basis vectors \(a_{i}\) .

